\section{Reference Validation and Deterministic Replay}

The repository contains three complementary validation layers. JSON Schema checks structure and required fields. A semantic validator checks invariants that cannot be expressed safely as simple types, including consistency between terminal states, timing eligibility, recovery references, missingness, and censoring. Deterministic replay tests reconstruct controller decisions from versioned telemetry and configuration.

\paragraph{Predeclared conformance cases.}
The machine-readable corpus manifest specifies three well-formed reference events: observed recovery, right-censored recovery, and an event in which a changed measurement context is explicitly marked as invalid for recovery inference. All three are expected to be accepted as \emph{representations}; accepting the third does not make its recovery measurement scientifically valid. The generator also defines seven deterministic negative mutations targeting a missing required schema version, mutually incompatible observed/censored terminal states, failed challenge eligibility, a mismatched reference rule, an out-of-domain recovery score, inconsistent latency arithmetic, and an undersized timing reference. These are expected to be rejected. The case inventory is a \emph{prespecified test design}, not an observed pass count.

\paragraph{Evaluation and failure reporting.}
The conformance generator emits per-case JSON and CSV records with case identifier, class, expected and observed acceptance, target invariant, outcome, and validator detail. Its own \texttt{pass} flag checks acceptance versus rejection; the target-invariant label alone is not proof that a negative case failed for the intended reason. A separate regression test, \nolinkurl{tests/test_s02_conformance_diagnostics.py}, executes the generator in a temporary directory and checks the declared invariant together with a distinguishing diagnostic location and reason for each of the seven negative mutations. This guards against counting a rejection caused by an unrelated invariant as evidence of the intended check. Diagnostic strings are part of this regression assertion; an intentional change to validator wording requires inspection rather than weakening the expected reason.

\paragraph{Prespecified acceptance rule.}
Before inspecting generated outputs, the case-level acceptance criterion requires acceptance of all three reference representations, rejection of all seven negative mutations, and the intended diagnostic location and reason for every negative rejection. Report separately the reference acceptance fraction $a_{\mathrm{ref}}/3$, negative rejection fraction $r_{\mathrm{neg}}/7$, and diagnostic-specific rejection fraction $d_{\mathrm{neg}}/7$. These denominators describe the designed corpus, not participants or observed performance. A rejected reference event, accepted negative mutation, or rejection for an unrelated reason fails its respective gate even if a coarser aggregate appears favorable. Preserve the complete case-level output, including failures; do not tune the corpus, expected diagnostics, or validator against these outputs without recording a new revision and rerunning the entire gate. Replay agreement and queue idempotence have separate evidence requirements and cannot be inferred from these fractions.

\paragraph{Reproducible execution boundary.}
At a recorded commit, execute the diagnostic-specificity unittest, the schema-table drift check, the conformance generator, and the independent controller-replay and queue tests. Archive the actual unittest output and both generated conformance files (\nolinkurl{conformance_report.json} and \nolinkurl{conformance_cases.csv}) with the schema blob, commit SHA, command, Python version, and execution environment. A passing generator summary is insufficient to claim diagnostic specificity unless the separate regression test also passes. Report every accepted negative mutation, rejected reference event, unexpected diagnostic, replay mismatch, or queue failure. The existence of these scripts and tests is not evidence that they have been executed at the manuscript revision; no numerical conformance result is asserted here. The validator must also reject non-finite numeric inputs, whether supplied directly to the semantic validation interface or introduced through nonstandard JSON tokens or overflowing numeric literals. Such values must not be allowed to bypass threshold and model-consistency checks.

\paragraph{Evidence traceability and failure classification.}
For each validation layer, the archived evidence record must identify the exact source revision, input fixture or event-sequence identifier, expected behavior, observed behavior, executable command, dependency versions, and machine-readable outcome. Classify a gate as \texttt{PASS} only when its prescribed assertions execute successfully at the recorded revision; use \texttt{FAIL} for an executed assertion that fails, \texttt{NOT\_RUN} when execution has not occurred, and \texttt{BLOCKED} when an identified environmental or access constraint prevents execution. Preserve failing records and stderr alongside successful records. Schema conformance, semantic diagnostics, controller replay, and queue idempotence require separate gate records, so an aggregate successful command cannot conceal an unexecuted layer. Record any post hoc change to a fixture, invariant, diagnostic expectation, or test command as a new evidence revision rather than silently replacing the prespecified corpus.

\paragraph{Independent validation layers.}
Conformance of individual event objects must not be conflated with deterministic controller replay or with idempotent queue/acknowledgement behavior. Replay requires comparison with versioned expected controller decisions, while queue testing requires separate duplicate-delivery and acknowledgement scenarios. The primary methods-paper outcomes are software-level: reference-event acceptance, negative-case rejection, replay agreement, queue idempotence, and preservation of missingness and right-censoring semantics. Their numerical results remain unreported until the corresponding tests have been executed and archived at a specified revision.

These tests can establish internal auditability and consistency of the software contract. They cannot establish the clinical validity of the recovery construct, transfer to other games, or intervention efficacy.
